% GENERATED FILE. Edit canonical lessons, then run npm run generate:books.
% canonical-source-hash: fnv1a64:14202db9e786d1ab
% canonical-lessons: PT-C33-pele, PT-C33-unha, PT-C33-queixo, PT-C33-garganta, PT-C33-estomago

\chapter{Skin, Nail, Chin, Throat, Stomach}
\label{ch:pt-skin-nail-chin-throat-stomach}
\hlchaptermodality{33}

\begin{chapteropening}
\textbf{By the end of this chapter:} \emph{I can name skin, a nail, the chin, the throat and the stomach.}

Complete the last lesson of chapter 33: I can name skin, a nail, the chin, the throat and the stomach.
\end{chapteropening}

\section[a pele]{a pele — skin}

\label{lesson:PT-C33-pele}



\begin{quote}
Before the new one: say the Portuguese for the face, then the Portuguese for blood.
\end{quote}

\subsection*{You'll want to know: a pele}
\textbf{a pele} — \textquotedblleft{}skin\textquotedblright{}.

\begin{cousinweb}
\textbf{pele} $\leftarrow$ Latin \emph{pellis}, \textquotedblleft{}skin, hide\textquotedblright{} (EtymDB 2.1) — inherited.
\end{cousinweb}

\begin{grammarlens}[title={What you've built}]
One more word for this chapter.
\end{grammarlens}

\subsection*{Your turn}
\begin{itemize}
\raggedright
  \item \emph{Say it:} \emph{a pele}
  \item \emph{Say it:} \emph{a pele}, once more
  \item \emph{Say it:} say \emph{a pele}
  \item \emph{Recall:} say the Portuguese for the face, then the Portuguese for blood, then say \emph{a pele} again
\end{itemize}

\begin{tcolorbox}[breakable,colback=teal!4,colframe=teal!35!black,title={Before you move on}]
What does a pele mean? (\textquotedblleft{}Skin\textquotedblright{}.) Say it once more.
\end{tcolorbox}
\section[a unha]{a unha — a nail (finger or toe)}

\label{lesson:PT-C33-unha}



\begin{quote}
Before the new one: say the Portuguese for blood, then the Portuguese for skin.
\end{quote}

\subsection*{You'll want to know: a unha}
\textbf{a unha} — \textquotedblleft{}a nail (finger or toe)\textquotedblright{}.

\begin{cousinweb}
\textbf{unha} $\leftarrow$ Latin \emph{ungula}, \textquotedblleft{}claw, hoof\textquotedblright{} (EtymDB 2.1), through Vulgar Latin \emph{*ungla} — inherited.
\end{cousinweb}

\begin{grammarlens}[title={What you've built}]
One more word for this chapter.
\end{grammarlens}

\subsection*{Your turn}
\begin{itemize}
\raggedright
  \item \emph{Say it:} \emph{a unha}
  \item \emph{Say it:} \emph{a unha}, once more
  \item \emph{Say it:} say \emph{a unha}
  \item \emph{Recall:} say the Portuguese for blood, then the Portuguese for skin, then say \emph{a unha} again
\end{itemize}

\begin{tcolorbox}[breakable,colback=teal!4,colframe=teal!35!black,title={Before you move on}]
What does a unha mean? (\textquotedblleft{}A nail (finger or toe)\textquotedblright{}.) Say it once more.
\end{tcolorbox}
\section[o queixo]{o queixo — the chin}

\label{lesson:PT-C33-queixo}



\begin{quote}
Before the new one: say the Portuguese for skin, then the Portuguese for a nail.
\end{quote}

\subsection*{You'll want to know: o queixo}
\textbf{o queixo} — \textquotedblleft{}the chin\textquotedblright{}.

\begin{cousinweb}
\textbf{queixo} $\leftarrow$ Vulgar Latin \emph{*capseum} (EtymDB 2.1), from Latin \emph{capsa}, \textquotedblleft{}box\textquotedblright{} — inherited.
\end{cousinweb}

\begin{grammarlens}[title={What you've built}]
One more word for this chapter.
\end{grammarlens}

\subsection*{Your turn}
\begin{itemize}
\raggedright
  \item \emph{Say it:} \emph{o queixo}
  \item \emph{Say it:} \emph{o queixo}, once more
  \item \emph{Say it:} say \emph{o queixo}
  \item \emph{Recall:} say the Portuguese for skin, then the Portuguese for a nail, then say \emph{o queixo} again
\end{itemize}

\begin{tcolorbox}[breakable,colback=teal!4,colframe=teal!35!black,title={Before you move on}]
What does o queixo mean? (\textquotedblleft{}The chin\textquotedblright{}.) Say it once more.
\end{tcolorbox}
\section[a garganta]{a garganta — the throat}

\label{lesson:PT-C33-garganta}



\begin{quote}
Before the new one: say the Portuguese for a nail, then the Portuguese for the chin.
\end{quote}

\subsection*{You'll want to know: a garganta}
\textbf{a garganta} — \textquotedblleft{}the throat\textquotedblright{}.

\begin{cousinweb}
\textbf{garganta} is built on an imitative root \emph{garg-}, the sound of the throat; English \emph{gargle} comes from the same sound. Its exact formation is uncertain.
\end{cousinweb}

\begin{grammarlens}[title={What you've built}]
One more word for this chapter.
\end{grammarlens}

\subsection*{Your turn}
\begin{itemize}
\raggedright
  \item \emph{Say it:} \emph{a garganta}
  \item \emph{Say it:} \emph{a garganta}, once more
  \item \emph{Say it:} say \emph{a garganta}
  \item \emph{Recall:} say the Portuguese for a nail, then the Portuguese for the chin, then say \emph{a garganta} again
\end{itemize}

\begin{tcolorbox}[breakable,colback=teal!4,colframe=teal!35!black,title={Before you move on}]
What does a garganta mean? (\textquotedblleft{}The throat\textquotedblright{}.) Say it once more.
\end{tcolorbox}
\section[o estômago]{o estômago — the stomach}

\label{lesson:PT-C33-estomago}



\begin{quote}
Before the new one: say the Portuguese for the chin, then the Portuguese for the throat.
\end{quote}

\subsection*{You'll want to know: o estômago}
\textbf{o estômago} — \textquotedblleft{}the stomach\textquotedblright{}.

\begin{cousinweb}
\textbf{estômago} $\leftarrow$ Latin \emph{stomachus}, \textquotedblleft{}stomach\textquotedblright{} (EtymDB 2.1) — a learned borrowing from Latin. Latin took it from Greek \emph{stomachos}, from \emph{stoma}, \textquotedblleft{}mouth\textquotedblright{}.
\end{cousinweb}

\begin{grammarlens}[title={What you've built}]
One more word for this chapter.
\end{grammarlens}

\subsection*{Your turn}
\begin{itemize}
\raggedright
  \item \emph{Say it:} \emph{o estômago}
  \item \emph{Say it:} \emph{o estômago}, once more
  \item \emph{Say it:} say \emph{o estômago}
  \item \emph{Recall:} say the Portuguese for the chin, then the Portuguese for the throat, then say \emph{o estômago} again
\end{itemize}

\begin{tcolorbox}[breakable,colback=teal!4,colframe=teal!35!black,title={Before you move on}]
What does o estômago mean? (\textquotedblleft{}The stomach\textquotedblright{}.) Say it once more.
\end{tcolorbox}
